\label{chap.3}
\section{Introduction}
Different kinds of data are required at each stage of the pipeline and no one source contains all of these. This chapter outlines the datasets gathered for the project, the reasons that none of these constitute a single dataset and the limitations associated with each. Table~\ref{tab:datasets} lists the datasets that were collected and used.

\section{Datasets Used}
Table~\ref{tab:datasets} gives the size, contents, role and preprocessing of each dataset.

{\footnotesize
\begin{longtable}{|p{2.3cm}|p{2.5cm}|p{3.3cm}|p{2.6cm}|p{2.6cm}|}
\caption{Datasets used in the project}
\label{tab:datasets}\\
\hline
\textbf{Dataset} & \textbf{Size} & \textbf{Contents} & \textbf{Role} & \textbf{Preprocessing} \\ \hline
\endfirsthead
\hline
\textbf{Dataset} & \textbf{Size} & \textbf{Contents} & \textbf{Role} & \textbf{Preprocessing} \\ \hline
\endhead
Telecom Italia Milan & 62 days, 10,000 cells, 10-minute intervals & SMS, call and internet activity on a $100 \times 100$ grid of 235\,m cells & Real traffic for SSL, TGNN, the unified environment and the Digital Twin dataset & Internet, calls and SMS mapped to eMBB, URLLC and mMTC demand; per-cell 95th-percentile scaling for the city-wide forecaster \\ \hline
5G-NIDD & 1{,}215{,}890 flows (408{,}655 unique) & Labelled benign and attack flows from a 5G network & Training and evaluation of the SSL encoder & Exact duplicates and identifier columns removed{,} log scaling{,} one-hot encoding{,} stratified 80/20 split \\ \hline
NeversNet5G & 716 UE files, about 27\,GB, 19 gNodeBs & Vehicular 5G NR telemetry produced with Simu5G and SUMO & UE to gNodeB topology (695 nodes) and telemetry replay & Graph built by nearest-gNodeB association \\ \hline
B5G network slicing & 7,984 samples & Transport topologies with eMBB, URLLC and mMTC flows & First MARL environment & Train, validation and test split \\ \hline
6G RAN telemetry & 30,000 windows & PRB load, latency percentiles, SLA flag and reward & Calibration of latency against load & Latency and jitter curves fitted against load \\ \hline
DeepSlice, Secure5G, Kaggle slicing sets, Botez 6G slicing & Small & Slice labels and QoS budgets & Slice-type benchmarks & --- \\ \hline
\end{longtable}
}

\section{Is there a Unified Public Dataset?}
A dataset that links all the stages of the pipeline would have to hold, in one record, the topology of the network, the traffic of each cell over time, the demand of each slice, QoS measurements, the SLA outcome, and the allocation action taken with its reward. We searched the datasets we hold and the web, and none of them contains all of these.

The closest is the B5G slicing dataset, which covers only the transport network and has no record of allocation actions. Published O-RAN slicing traces come from single-cell testbeds, and we could not confirm that they are released. The literature audit of Chapter~\ref{chap.2} agrees: of the sixteen papers reviewed, eight use their own data sets, and none releases a public dataset with topology, per-slice traffic, QoS and allocation decisions together.

On the advice of our guide we therefore built a unified dataset synthetically, based on real data wherever real data exist. The environment is described in Chapter~\ref{chap.4} and the data it generates in Chapter~\ref{chap.8}.

\section{Dataset Limitations}
Each source limits what the experiments can show.
\begin{itemize}
	\item \textbf{Telecom Italia Milan} has no network-control information: no PRB allocation, slice configuration, latency, QoS or SLA decisions, and no MARL actions. It supplies the real traffic-demand layer only, and the demand of each slice is assumed from it.
	\item \textbf{B5G} models the transport network and not the RAN. The samples used do not include latency or loss, so the reward of the first MARL environment is nearly trivial, and no Digital Twin loop could approve actions on it.
	\item \textbf{6G RAN telemetry} calibrates latency against load, but the telemetry is itself simulated. The calibration therefore ties the environment to a model and not to measurements from an operating network.
	\item \textbf{NeversNet5G} was produced with Simu5G and SUMO and not on a real deployment.
	\item \textbf{5G-NIDD} is a labelled intrusion-detection dataset. It can evaluate the SSL encoder, but it holds no slicing information. About two thirds of its flows are exact duplicates (807{,}235 of 1{,}215{,}890), which would leak between a random training and test split, so only the unique flows are used.
\end{itemize}
